\documentclass[../../main.tex]{subfiles}
\begin{document}

{\bf [解]}

与えられた漸化式
\begin{align}
 \begin{cases}
   a_1=\dfrac{1}{n(n+1)} \\  
   a_{k+1}=-\dfrac{1}{k+n+1}+\dfrac{n}{k}\sum_{i=1}^ka_i\qquad(k=1,2,\dots).\tag{①}
 \end{cases} \label{eq:1}
\end{align}
について考える．

\paragraph{(1)}

\cref{eq:1}に$k=1,2$を順に代入すると
\begin{align}
a_2
  &=-\dfrac{1}{n+2}+n\cdot\dfrac{1}{n(n+1)} \\
  &=-\dfrac{1}{n+2}+\dfrac{1}{n+1} \\
  &=\dfrac{1}{(n+1)(n+2)}
\end{align}
および
\begin{align}
 a_3
  &=-\dfrac{1}{n+3}+\dfrac{n}{2}\Bigl(\dfrac{1}{n(n+1)}+\dfrac{1}{(n+1)(n+2)}\Bigr) \\
  &=-\dfrac{1}{n+3}+\dfrac{n}{2}\cdot\dfrac{1}{n+1}\cdot\dfrac{2(n+1)}{n(n+2)} \\
  &=-\dfrac{1}{n+3}+\dfrac{1}{n+2}=\dfrac{1}{(n+2)(n+3)}
\end{align}
を得る．

\paragraph{(2)}

数学的帰納法により
\begin{align}
 a_k=\dfrac{1}{(n+k-1)(n+k)} \label{eq:2}
\end{align}
であることを示す．
$k=1$のときは(1)より$a_1=\dfrac{1}{n(n+1)}$と一致して成立する．
そこで$k\le m$（$m\in\mathbb N$）なる全ての$k$で成立すると仮定する．すると$a_i$の和
\begin{align}
 S_m=\displaystyle\sum_{i=1}^ma_i
\end{align}
は
\begin{align}
 S_m
  &=\sum_{i=1}^m\Bigl(\dfrac{1}{n+i-1}-\dfrac{1}{n+i}\Bigr) \\
  &=\dfrac{1}{n}-\dfrac{1}{n+m}
\end{align}
と計算できるから，\cref{eq:1}の漸化式に代入して$a_{m+1}$は
\begin{align}
 a_{m+1}
  &=-\dfrac{1}{n+m+1}+\dfrac{n}{m}\Bigl(\dfrac{1}{n}-\dfrac{1}{n+m}\Bigr) \\
  &=-\dfrac{1}{n+m+1}+\dfrac{n}{m}\cdot\dfrac{m}{n(n+m)} \\
  &=-\dfrac{1}{n+m+1}+\dfrac{1}{n+m} \\
  &=\dfrac{1}{(n+m)(n+m+1)}
\end{align}
となり，これは$k=m+1$でも\cref{eq:2}が成立することを示す．よって数学的帰納法により，すべての$k$で
\begin{align}
a_k=\dfrac{1}{(n+k-1)(n+k)}
\end{align}
である．

\paragraph{(3)}

\begin{align}
 b_n=\displaystyle\sum_{k=1}^n\sqrt{a_k} \label{eq:3}
\end{align}
について考える．

自明な関係式$(n+k-1)^2<(n+k-1)(n+k)<(n+k)^2$の逆数をとって，\cref{eq:2}より
\begin{align}
 &\dfrac{1}{(n+k)^2}<a_k<\dfrac{1}{(n+k-1)^2} \\
\therefore\ 
 &\dfrac{1}{n+k}<\sqrt{a_k}<\dfrac{1}{n+k-1} \label{eq:4}
\end{align}
だから，$k=1,2,\cdots,n$について和をとって\cref{eq:3}を代入すると
\begin{align}
 \sum_{k=1}^n\dfrac{1}{n+k}<b_n<\sum_{k=1}^n\dfrac{1}{n+k-1} \label{eq:5}
\end{align}
である．

さて，両端の和は$n\to\infty$とすると区分求積法により
\begin{align}
\sum_{k=1}^n\dfrac{1}{n+k} 
  &=\dfrac{1}{n}\sum_{k=1}^n\dfrac{1}{1+\dfrac{k}{n}} \\
  &\xrightarrow[n\to\infty]{}\int_0^1\dfrac{dx}{1+x} =\log2
\end{align}
\begin{align}
 \sum_{k=1}^n\dfrac{1}{n+k-1}
  &=\dfrac{1}{n}\sum_{k=1}^n\dfrac{1}{1+\dfrac{k-1}{n}} \\
  &\xrightarrow[n\to\infty]{}\int_0^1\dfrac{dx}{1+x} =\log2
\end{align}
だから，両方$\log 2$に収束する．従って\cref{eq:5}において挟み撃ちの原理より
\begin{align}
\lim_{n\to\infty}b_n=\log2
\end{align}
である．



\newpage

\end{document}
